%!TEX TS-program = pdflatex
%!TEX encoding = UTF-8 Unicode
% Jake's Resume, a Quarto format for the LaTeX template packaged at
% https://github.com/sb2nov/resume by Sourabh Bajaj -- the single most-copied
% resume template on Overleaf, where it circulates under the nickname "Jake's
% Resume".
%
% Adapted to be a Quarto template by Mitchell O'Hara-Wild
%
% Upstream license: MIT (see LICENSE.md)
%
% Unlike Awesome-CV or latexcv, upstream ships no class file: it is one
% ~200-line .tex of \documentclass{article} plus a handful of \newcommand
% definitions. Everything up to \begin{document} -- those commands, class
% options, colours and Pandoc's own preamble -- lives here; sb2nov.tex adds
% only the document-specific header, since upstream hard-codes its name and
% contact details rather than defining a command for them.
%-------------------------------------------------------------------------------
% CONFIGURATIONS
%-------------------------------------------------------------------------------
$passoptions.latex()$
%
% US letter by default, matching upstream; set `papersize: a4` to override.
$doc-class.tex()$

\usepackage{titlesec}
% Upstream loads the plain `color` package, which has no notion of `headcolor`
% as a hex string. `xcolor` is loaded instead -- with the same `usenames,
% dvipsnames` options upstream passes, which it also understands -- since only
% xcolor's `HTML` colour model and \colorlet can parse and fall back on one,
% exactly as awesome-cv and latexcv already do.
\usepackage[usenames,dvipsnames]{xcolor}
\usepackage{enumitem}
\usepackage{fancyhdr}
\usepackage{geometry}
% Provides \urlstyle below. Upstream gets this for free as a side effect of
% loading hyperref itself; hyperref is loaded later here, by common.latex(),
% so \url's own package is loaded directly rather than moving \urlstyle{same}
% out of its upstream position.
\usepackage{url}
% Upstream also loads `tabularx`, but never opens the `tabularx` environment
% it provides -- every table in the file is a plain `tabular*`. Dropped as
% dead weight rather than carried over unused.

% This template writes its own preamble rather than deferring to Pandoc's
% default article one, so nothing sets input or font encoding unless it is
% done here. `iftex` guards both: they are pdftex-only, and a document that
% overrides `pdf-engine` to xelatex or lualatex gets UTF-8 input and modern
% font handling natively, without these.
\usepackage{iftex}
\ifPDFTeX
  \usepackage[utf8]{inputenc}
  \usepackage[$if(fontenc)$$fontenc$$else$T1$endif$]{fontenc}
  % Ensures the generated PDF's text is machine-readable/ATS-parsable --
  % the entire point of this design -- by teaching pdfTeX the Unicode glyph
  % each ligature and accent command actually stands for. Guarded the same
  % way: the primitive it sets does not exist outside pdfTeX.
  \input{glyphtounicode}
  \pdfgentounicode=1
\fi

% Upstream loads `babel[english]` itself. Pandoc's own babel-lang.tex, pulled
% in by common.latex() below, already loads babel from the document's `lang`
% metadata (English unless told otherwise), so loading it a second time here
% would risk an option clash instead of adding anything.

\pagestyle{fancy}
\fancyhf{} % Clear all header and footer fields
\fancyfoot{}
\renewcommand{\headrulewidth}{0pt}
\renewcommand{\footrulewidth}{0pt}

% Page margins. Upstream reaches its (approximately half-inch) margins with
% `\usepackage[empty]{fullpage}` plus a stack of `\addtolength` corrections on
% top of it, which only add up for letterpaper -- change the paper size and
% the corrections no longer land where upstream intended. `geometry` is used
% instead, as every other quarto-vitae template does, so `papersize` and
% `geometry` in the document front matter both work as expected.
$if(geometry)$
\geometry{$for(geometry)$$geometry$$sep$,$endfor$}
$else$
\geometry{left=0.75in, right=0.75in, top=0.6in, bottom=0.7in}
$endif$

\urlstyle{same}

\raggedbottom
\raggedright
\setlength{\tabcolsep}{0in}

% Accent colour for the section rule and heading text.
%
% Jake's Resume itself is black-and-white -- its one \color{black} call in
% \titleformat below is there only to stop \color{Blue} et al from a
% surrounding document leaking into the section rule, not to make the colour
% configurable. `headcolor` is honoured anyway, for consistency with every
% other quarto-vitae template, but defaults to upstream's own black so a
% document that says nothing about it renders identically to upstream.
$if(headcolor)$
\definecolor{sectioncolor}{HTML}{$headcolor$}
$else$
\colorlet{sectioncolor}{black}
$endif$

% Sections formatting
\titleformat{\section}{
  \vspace{-4pt}\scshape\raggedright\large\color{sectioncolor}
}{}{0em}{}[\color{sectioncolor}\titlerule \vspace{-5pt}]

% Upstream defines no second-level heading at all -- every upstream section is
% a \resumeSubheading table row instead. A `##` is still given a styled
% command rather than falling through to article's plain, unindented
% \subsection, so a document that does use one degrades gracefully.
\titleformat{\subsection}{
  \vspace{-2pt}\scshape\raggedright\normalsize\color{sectioncolor}
}{}{0em}{}

%-------------------------------------------------------------------------------
% Custom commands
%
% Kept verbatim from upstream, so the entries/ templates below call exactly
% the commands sb2nov/resume itself calls; only how those commands are
% invoked (which what/with/where/when/why field goes in which argument) is
% this port's own work. See entries/detailed.tex, entries/brief.tex and
% entries/skills.tex for the wiring.
\newcommand{\resumeItem}[2]{
  \item\small{
    \textbf{#1}{: #2 \vspace{-2pt}}
  }
}

% Just in case someone needs a heading that does not need to be in a list
\newcommand{\resumeHeading}[4]{
    \begin{tabular*}{0.99\textwidth}[t]{l@{\extracolsep{\fill}}r}
      \textbf{#1} & #2 \\
      \textit{\small #3} & \textit{\small #4} \\
    \end{tabular*}\vspace{-5pt}
}

\newcommand{\resumeSubheading}[4]{
  \vspace{-1pt}\item
    \begin{tabular*}{0.97\textwidth}[t]{l@{\extracolsep{\fill}}r}
      \textbf{#1} & #2 \\
      \textit{\small #3} & \textit{\small #4} \\
    \end{tabular*}\vspace{-5pt}
}

\newcommand{\resumeSubSubheading}[2]{
    \begin{tabular*}{0.97\textwidth}{l@{\extracolsep{\fill}}r}
      \textit{\small #1} & \textit{\small #2} \\
    \end{tabular*}\vspace{-5pt}
}

\newcommand{\resumeItemPlain}[1]{
  \item\small{
    {#1 \vspace{-2pt}}
  }
}

\newcommand{\resumeSubItem}[2]{\resumeItem{#1}{#2}\vspace{-4pt}}

\renewcommand{\labelitemii}{$$\circ$$}

\newcommand{\resumeSubHeadingListStart}{\begin{itemize}[leftmargin=*]}
\newcommand{\resumeSubHeadingListEnd}{\end{itemize}}
\newcommand{\resumeItemListStart}{\begin{itemize}}
\newcommand{\resumeItemListEnd}{\end{itemize}\vspace{-5pt}}

%-------------------------------------------------------------------------------
%	PERSONAL INFORMATION
%
%	Upstream hard-codes its own name, one link and two contact lines directly
%	in the header table. This template needs that stood up from the document's
%	front matter instead, and with more than two contact fields possible, so
%	the right-hand column is built from \jrcontact calls that separate
%	themselves, in the manner of latexcv's \cvcontacts -- a fixed `\\` between
%	every conditional block would otherwise leave a blank stacked line wherever
%	a field is unset.
%-------------------------------------------------------------------------------
\newif\ifjrfirstcontact
\newcommand{\jrcontactsep}{\\}
\newcommand{\jrcontact}[1]{\ifjrfirstcontact\jrfirstcontactfalse\else\jrcontactsep\fi #1}
% Every line ends in a `%', as latexcv's own \cvcontacts does: an unset field
% leaves its whole line blank, and an unescaped blank line is a paragraph
% break. Left in, that \par lands *between* two \jrcontact calls and closes
% the paragraph \jrcontactsep's `\\` needs still open, so the next contact
% after a gap fails with "There's no line here to end" instead of stacking.
\newcommand{\jrcontacts}{%
  \jrfirstcontacttrue
  $if(email)$\jrcontact{Email: \href{mailto:$email$}{$email$}}$endif$%
  $if(phone)$\jrcontact{Mobile: \href{tel:$phone$}{$phone$}}$endif$%
  $if(www)$\jrcontact{\href{https://$www$}{$www$}}$endif$%
  $if(github)$\jrcontact{\href{https://github.com/$github$}{github.com/$github$}}$endif$%
  $if(linkedin)$\jrcontact{\href{https://www.linkedin.com/in/$linkedin$}{linkedin.com/in/$linkedin$}}$endif$%
}

% Computer Modern's default fonts are Metafont bitmaps, not scalable outlines.
% common.latex() below turns on microtype's font expansion unconditionally,
% which errors outright the moment it meets a non-scalable font. lmodern is
% Computer Modern's metric-compatible, fully scalable replacement, and is
% loaded before common.latex() so that microtype's own font detection sees it
% -- loaded by default here for the same reason latexcv defaults to raleway,
% except that upstream sets no font of its own to take that default's place.
% A document naming its own `fontfamily` is trusted to have made expansion
% work, or to not mind microtype falling back silently if it has not.
$if(fontfamily)$
$else$
\ifPDFTeX
  \usepackage{lmodern}
\fi
$endif$

%-------------------------------------------------------------------------------
%	PANDOC SUPPORT
%-------------------------------------------------------------------------------
% Upstream's own \usepackage[hidelinks]{hyperref} is not written here: the
% common.latex() partial below loads hyperref itself (via pandoc.tex), and
% loading it twice with different options is a package clash. Passing the
% option through beforehand gets the same plain, unboxed, uncoloured links.
\PassOptionsToPackage{hidelinks}{hyperref}

% common.latex is required in every custom Quarto LaTeX template, and calls
% pandoc.tex, which in turn brings in syntax highlighting, tables.tex,
% graphics.tex, citations.tex, babel-lang.tex, tightlist.tex, biblio-config.tex
% and header-includes. Do not hand-copy any of those: they are Pandoc's to
% maintain, and a stale local copy silently drifts out of step with the reader.
$common.latex()$
$after-header-includes.latex()$
$hypersetup.latex()$

% Honours mainfont/sansfont/monofont under a unicode engine, and `fontfamily`
% (a LaTeX font package name, e.g. `mathpazo`) under pdflatex -- which is also
% how a document overrides the lmodern default set above. fonts.latex is NOT
% called: on top of loading lmodern unconditionally, doubling the one above,
% it would also load unicode-math, which pdflatex cannot process.
$font-settings.latex()$
